\documentclass[11pt,a4paper]{article}

% Tectonic = moteur XeTeX : UTF-8 natif
\usepackage[french]{babel}
\usepackage[a4paper, top=2.2cm, bottom=2.0cm, left=2.2cm, right=2.2cm]{geometry}
\usepackage{xcolor}
\usepackage{titlesec}
\usepackage{enumitem}
\usepackage{booktabs}
\usepackage{tcolorbox}
\usepackage{listings}
\usepackage{fancyhdr}

\definecolor{accent}{HTML}{1F4E79}
\definecolor{accent2}{HTML}{2E7D32}
\definecolor{soft}{HTML}{F2F5F9}
\definecolor{codebg}{HTML}{F6F8FA}
\definecolor{warn}{HTML}{B71C1C}

\usepackage[colorlinks=true, linkcolor=accent, urlcolor=accent,
    pdftitle={Guide - Creer le depot GitHub pour le projet de these},
    pdfauthor={Jihan El-khattabi},
    pdfsubject={Guide pas-a-pas - creation d'un depot GitHub},
    pdfkeywords={GitHub, Git, depot, push, commit, these}
]{hyperref}

\pagestyle{fancy}
\fancyhf{}
\setlength{\headheight}{14pt}
\fancyhead[L]{\small\color{accent}\textbf{Guide -- Créer le dépôt GitHub}}
\fancyhead[R]{\small\color{accent}Jihan El-khattabi -- LBECS}
\fancyfoot[C]{\small\thepage\,/\,\pageref{LastPage}}
\renewcommand{\headrulewidth}{0.4pt}
\renewcommand{\headrule}{\color{accent}\hrule width\headwidth height\headrulewidth}

\titleformat{\section}{\color{accent}\large\bfseries}{\thesection.}{0.5em}{}
\titlespacing*{\section}{0pt}{10pt}{4pt}
\setlength{\parskip}{4pt}
\setlength{\parindent}{0pt}
\setlength{\emergencystretch}{2em}

% Style des blocs de commandes
\lstdefinestyle{shell}{
  basicstyle=\ttfamily\small,
  backgroundcolor=\color{codebg},
  frame=single,
  rulecolor=\color{accent!60},
  framesep=5pt,
  framerule=0.4pt,
  breaklines=true,
  breakatwhitespace=true,
  showstringspaces=false,
  columns=fullflexible,
  keepspaces=true,
  keywordstyle=\color{accent}\bfseries,
  morekeywords={git,cd,sudo,ssh-keygen,cat,echo,ls,mkdir,ssh,git}
}
\lstset{style=shell}

% Boite d'etape
\newtcolorbox{etape}[1]{%
  colback=accent!8, colframe=accent, boxrule=0.5pt, arc=2pt,
  left=8pt, right=8pt, top=4pt, bottom=4pt,
  fonttitle=\bfseries\color{white}, colbacktitle=accent,
  title={#1}}
\newtcolorbox{astuce}{colback=accent2!8, colframe=accent2, boxrule=0.4pt,
  arc=2pt, left=6pt, right=6pt, top=3pt, bottom=3pt,
  title={\textbf{Astuce}}, fonttitle=\color{white}\bfseries, colbacktitle=accent2}
\newtcolorbox{attention}{colback=warn!6, colframe=warn, boxrule=0.4pt,
  arc=2pt, left=6pt, right=6pt, top=3pt, bottom=3pt,
  title={\textbf{Attention}}, fonttitle=\color{white}\bfseries, colbacktitle=warn}

\begin{document}

\begin{center}
{\color{accent}\Large\bfseries Guide pas à pas -- Créer le dépôt GitHub du projet de thèse}\\[3pt]
{\large\itshape Étude phytochimique de \textit{Rugulopteryx okamurae} dans la Méditerranée marocaine}\\[4pt]
\rule{0.97\linewidth}{0.6pt}\\[3pt]
\small
\textbf{Doctorante :} Jihan El-khattabi \quad|\quad
\textbf{ORCID :} \href{https://orcid.org/0009-0002-7946-4934}{0009-0002-7946-4934}\\
\textbf{Dépôt cible :} \texttt{github.com/JihanElkhattabi/rugulopteryx-okamurae-these}\\
\textbf{Date :} 15 novembre 2026
\end{center}
\vspace{2pt}

\section{Objectif et principe}
Ce guide vous accompagne pour mettre en ligne, sur GitHub, le dépôt local
\texttt{depot\_git/} déjà prêt (README, LICENSE, .gitignore, 7 commits, fichiers
d'exemple). À la fin, votre projet sera visible à l'adresse
\url{https://github.com/JihanElkhattabi/rugulopteryx-okamurae-these} et le lien
figurant dans votre rapport deviendra actif.

Le principe est simple : GitHub héberge une copie distante (\emph{remote}) de votre
dépôt local. Vous reliez les deux avec \texttt{git remote}, puis vous «~poussez~»
(\texttt{git push}) vos commits vers GitHub. Ensuite, chaque modification locale est
synchronisée par \texttt{add} + \texttt{commit} + \texttt{push}.

\section{Prérequis}
\begin{itemize}
  \item Un \textbf{compte GitHub} : créez-le gratuitement sur \url{https://github.com/signup}
        (utilisez votre identifiant \texttt{JihanElkhattabi} ou un nom similaire).
  \item \textbf{Git} installé sur votre machine. Vérifiez dans un terminal :
\begin{lstlisting}
git --version
\end{lstlisting}
        Si Git n'est pas installé :
        \begin{itemize}
          \item \textbf{Linux (Debian/Ubuntu)} : \texttt{sudo apt install git}
          \item \textbf{macOS} : \texttt{brew install git} (ou installez Xcode Command Line Tools)
          \item \textbf{Windows} : téléchargez Git sur \url{https://git-scm.com/download/win}
        \end{itemize}
  \item Le dossier \textbf{\texttt{depot\_git/}} fourni (contenant déjà 7 commits).
\end{itemize}

\section{Étape 1 -- Configurer Git (une seule fois)}
Avant le premier envoi, indiquez à Git votre identité (servira à signer chaque commit).
Remplacez l'e-mail par celui associé à votre compte GitHub :

\begin{lstlisting}
git config --global user.name "Jihan El-khattabi"
git config --global user.email "votre.email@example.com"
git config --global init.defaultBranch main
git config --global core.autocrlf input   # Linux/macOS ; mettre "true" sous Windows
\end{lstlisting}

Vérifiez :
\begin{lstlisting}
git config --global --list
\end{lstlisting}

\section{Étape 2 -- Créer le dépôt vide sur GitHub}
\begin{enumerate}
  \item Connectez-vous sur \url{https://github.com}.
  \item Cliquez sur l'icône \textbf{+} (en haut à droite) $\rightarrow$ \textbf{New repository}.
  \item Remplissez le formulaire :
        \begin{itemize}
          \item \textbf{Repository name} : \texttt{rugulopteryx-okamurae-these}
          \item \textbf{Description} : \textit{Écosystème numérique de thèse -- Étude phytochimique de Rugulopteryx okamurae}
          \item \textbf{Visibility} : \textbf{Public} (pour que le lien du rapport soit consultable)
          \item \textbf{Laissez DECOCHÉES} les cases «~Add a README~», «~.gitignore~» et
                «~License~» : votre dépôt local les contient déjà.
        \end{itemize}
  \item Cliquez sur \textbf{Create repository}.
\end{enumerate}

\begin{attention}
Si vous cochez «~Initialize this repository with a README~», GitHub créera un commit
initial \emph{distinct} du vôtre, ce qui provoquera un conflit au premier \texttt{push}.
Laissez donc le dépôt \textbf{strictement vide}.
\end{attention}

GitHub affiche alors une page avec l'URL du dépôt. Copiez-la, par exemple :
\url{https://github.com/JihanElkhattabi/rugulopteryx-okamurae-these.git}

\section{Étape 3 -- Connecter le dépôt local au dépôt distant}
Ouvrez un terminal à l'intérieur du dossier \texttt{depot\_git/} fourni (qui est déjà
un dépôt Git avec 7 commits), puis reliez-le au dépôt GitHub :

\begin{lstlisting}
cd /chemin/vers/depot_git

# Relier le depot local au depot distant nomme "origin"
git remote add origin https://github.com/JihanElkhattabi/rugulopteryx-okamurae-these.git

# S'assurer que la branche principale s'appelle bien "main"
git branch -M main

# Verifier la connexion
git remote -v
\end{lstlisting}

Vous devez voir deux lignes \texttt{origin} (\texttt{fetch} et \texttt{push}) pointant
vers votre URL GitHub.

\section{Étape 4 -- S'authentifier auprès de GitHub}
GitHub n'accepte plus votre mot de passe pour les opérations Git. Deux méthodes :

\textbf{Méthode A -- HTTPS + jeton d'accès personnel (PAT), recommandée pour débuter :}
\begin{enumerate}
  \item Sur GitHub : \textbf{Settings} $\rightarrow$ \textbf{Developer settings}
        $\rightarrow$ \textbf{Personal access tokens} $\rightarrow$ \textbf{Tokens (classic)}
        $\rightarrow$ \textbf{Generate new token}.
  \item Donnez un nom, une expiration, et cochez le scope \texttt{repo} (accès complet
        aux dépôts).
  \item Copiez le jeton généré (il ne s'affichera plus jamais).
  \item Au moment du \texttt{git push}, entrez votre nom d'utilisateur GitHub et,
        comme \textbf{mot de passe}, collez le jeton.
\end{enumerate}

\begin{astuce}
Pour éviter de ressaisir le jeton à chaque \texttt{push}, mettez-le en cache avec
l'une des commandes ci-dessous.
\end{astuce}
\begin{lstlisting}
git config --global credential.helper store     # simple (jeton en clair)
# ou, plus sur :
git config --global credential.helper "cache --timeout=86400"
\end{lstlisting}

\textbf{Méthode B -- SSH (une fois configurée, plus de jeton à gérer) :}
\begin{lstlisting}
ssh-keygen -t ed25519 -C "votre.email@example.com"   # genere une cle
cat ~/.ssh/id_ed25519.pub                             # affiche la cle publique
\end{lstlisting}
Copiez la clé publique et ajoutez-la sur GitHub : \textbf{Settings}
$\rightarrow$ \textbf{SSH and GPG keys} $\rightarrow$ \textbf{New SSH key}. Utilisez
alors l'URL SSH et remplacez l'URL HTTPS via la commande :
\begin{lstlisting}
git remote set-url origin git@github.com:JihanElkhattabi/rugulopteryx-okamurae-these.git
\end{lstlisting}

\section{Étape 5 -- Pousser le dépôt vers GitHub (premier push)}
\begin{lstlisting}
git push -u origin main
\end{lstlisting}

\begin{itemize}
  \item L'option \texttt{-u} (ou \texttt{--set-upstream}) enregistre \texttt{origin/main}
        comme branche de suivi : les prochains \texttt{git push} se feront sans argument.
  \item Git envoie vos \textbf{7 commits} et tous les fichiers associés.
  \item Si l'authentification réussit, la sortie affiche des lignes
        \texttt{* [new branch] main -> main} et un récapitulatif des objets envoyés.
\end{itemize}

\section{Étape 6 -- Vérifier sur GitHub}
\begin{enumerate}
  \item Rechargez la page de votre dépôt sur GitHub.
  \item Vous devez voir : le \texttt{README.md} affiché en page d'accueil, tous les
        dossiers (\texttt{docs/}, \texttt{data/}, \texttt{scripts/}, \texttt{results/}),
        le fichier \texttt{LICENSE}, et \texttt{7 commits} dans l'historique.
  \item Cliquez sur \textbf{commits} : la liste de vos 7 commits doit apparaître.
  \item Cliquez sur un fichier \texttt{.pdf} dans \texttt{docs/} : GitHub permet de le
       télécharger (et le lien de votre rapport devient actif).
\end{enumerate}

\section{Étape 7 -- Workflow quotidien (modifier, valider, envoyer)}
Une fois le dépôt en ligne, voici la routine pour chaque évolution :

\begin{lstlisting}
cd /chemin/vers/depot_git

# 1. Voir ce qui a change
git status

# 2. Ajouter les fichiers modifies (ou tout)
git add <fichier>
# ou:  git add .

# 3. Creer un commit avec un message clair
git commit -m "Ajout du chromatogramme de l'echantillon Tetouan-001"

# 4. Envoyer sur GitHub
git push
\end{lstlisting}

\begin{astuce}
Un bon message de commit décrit le \emph{quoi} et le \emph{pourquoi} en une phrase
courte, au présent ou à l'infinitif. Ex.~: «~Ajout des spectres RMN 13C du lot
Tetouan~» plutôt que «~modifs~».
\end{astuce}

\section{Étape 8 -- Mettre à jour le lien dans le rapport (optionnel)}
Si votre identifiant GitHub réel diffère de \texttt{JihanElkhattabi}, remplacez-le dans
les trois fichiers suivants, puis recompilez le PDF :
\begin{lstlisting}
Rapport_Ecosysteme_Numerique_These.tex
README.md
CITATION.cff
\end{lstlisting}

\begin{lstlisting}
# dans le dossier contenant le .tex
tectonic Rapport_Ecosysteme_Numerique_These.tex
\end{lstlisting}

\section{Aller plus loin}
\begin{itemize}
  \item \textbf{Description et sujets} : sur GitHub, cliquez sur l'engrenage
        $\rightarrow$ ajoutez des \texttt{Topics} (\texttt{phytochemistry},
        \texttt{macroalgae}, \texttt{invasive-species}, \texttt{open-science}) pour la
        découvrabilité.
  \item \textbf{Release / version} : étiquetez un jalon avec
        \texttt{git tag -a v1.0 -m "Version initiale"} puis \texttt{git push origin v1.0}.
  \item \textbf{DOI persistant} : reliez votre dépôt à \href{https://zenodo.org}{Zenodo}
        (\textbf{Settings} GitHub $\rightarrow$ \textbf{Webhooks} $\rightarrow$ Zenodo).
        Chaque release générera un DOI citable, idéal pour une thèse.
  \item \textbf{GitHub Pages} : pour afficher le README comme une page web, activez
        \textbf{Settings} $\rightarrow$ \textbf{Pages} $\rightarrow$ branche \texttt{main}.
\end{itemize}

\section{Dépannage (erreurs fréquentes)}

\textbf{1.~«~remote origin already exists~».} L'origine est déjà définie. Mettez plutôt
à jour l'URL :
\begin{lstlisting}
fatal: remote origin already exists.

# Solution : remplacer l'URL du distant
git remote set-url origin https://github.com/JihanElkhattabi/rugulopteryx-okamurae-these.git
\end{lstlisting}

\textbf{2.~Erreur d'authentification.} GitHub n'accepte plus les mots de passe.
Utilisez un \textbf{PAT} (cf. étape~4) ou une clé \textbf{SSH}.
\begin{lstlisting}
fatal: Authentication failed
Support for password authentication was removed
\end{lstlisting}

\textbf{3.~Push rejeté (non-fast-forward).} Le dépôt distant contient des commits
absents localement. Récupérez-les puis poussez :
\begin{lstlisting}
! [rejected] main -> main (fetch first)

# Solution : recuperer puis pousser
git pull --rebase origin main
git push
\end{lstlisting}

\textbf{4.~Échec du push à cause d'un gros fichier.} GitHub limite les fichiers à
100~Mo. Ajoutez-les à \texttt{.gitignore} ou utilisez \textbf{Git LFS}
(\texttt{git lfs install}). Pour les données brutes volumineuses (spectres RMN),
déposez-les plutôt sur Zenodo/Nakala (cf. PGD) et ne versionnez que le code et les
tables réduites.
\begin{lstlisting}
error: failed to push some refs
\end{lstlisting}

\textbf{5.~Les accents ou caractères s'affichent mal sur GitHub.} C'est un problème
d'encodage. Assurez-vous que vos fichiers sont en UTF-8
(\texttt{file -i <fichier>} sous Linux). Les fichiers de ce dépôt sont déjà en UTF-8.

\section{Récapitulatif en une page}
\begin{etape}{Séquence complète (depuis le dossier depot\_git/ déjà initialisé)}
Toutes les commandes essentielles, à exécuter dans l'ordre ci-dessous :
\end{etape}
\begin{lstlisting}
# 1. Identite (une fois)
git config --global user.name "Jihan El-khattabi"
git config --global user.email "votre.email@example.com"

# 2. Creer le depot VIDE sur github.com (sans README, sans licence)

# 3. Relier et pousser
cd depot_git
git remote add origin https://github.com/JihanElkhattabi/rugulopteryx-okamurae-these.git
git branch -M main
git push -u origin main

# 4. Ensuite, a chaque modification :
git add .
git commit -m "description courte du changement"
git push
\end{lstlisting}

À l'issue de ces étapes, le dépôt est en ligne, le lien du rapport est actif, et la
condition «~au moins 3 commits~» de la Slide~38 est largement satisfaite (7 commits).

Pour pérenniser le projet, pensez à créer une \textbf{release} étiquetée
(\texttt{git tag -a v1.0 -m "..."} puis \texttt{git push origin v1.0}) à chaque jalon
majeur de la thèse, et à relier le dépôt à \href{https://zenodo.org}{Zenodo} pour
obtenir un \textbf{DOI} citable. Ce DOI pourra figurer dans votre profil ORCID et dans
les publications issues de la thèse, fermant la boucle de l'écosystème numérique.

\label{LastPage}
\end{document}
